\section{Vorbereitung der Versuchanordnung}
\subsection{Justage des NaI-Szintillationsdetektor}

Zunächst wird die Hochspannung des Photomultipliers so eingestellt, dass der Detektor in seinem stabilen Arbeitsbereich betrieben wird. Die Spannung wird schrittweise erhöht, bis eine ausreichende Impulsrate und eine deutlich erkennbare Impulsamplitude erreicht sind, ohne den Photomultiplier in seinen nichtlinearen Arbeitsbereich zu treiben. Das Ausgangssignal des Vorverstärkers wird anschließend direkt auf einem Oszilloskop betrachtet; dabei werden der charakteristische Signalverlauf, die Impulsdauer und die Amplitude aufgezeichnet. Der Signalverlauf wird bestimmt durch das Szintillationsereignis im NaI(Tl)-Kristall, die Lichtemission, die interne Verstärkung des Photomultipliers sowie die nachfolgende Impulsformung durch den Vorverstärker. Dieses Signal durchläuft anschließend eine weitere Verarbeitung im Hauptverstärker, bevor es zur digitalen Spektrenaufnahme an einen Vielkanalanalysator (MCA) weitergeleitet wird.\\

Im nächsten Schritt werden Gammaspektren der verwendeten radioaktiven Quellen aufgenommen. Konkret werden nacheinander die Spektren von $^{60}Co$, $^{137}Cs$ und $^{152}Eu$ aufgezeichnet. Zusätzlich wird das Untergrundspektrum (ohne Quelle) gemessen, um eine spätere Subtraktion der Untergrundbeiträge zu ermöglichen. Die Messzeiten werden so gewählt, dass die Photopeaks scharf ausgeprägt und statistisch gut aufgelöst sind.\\

Auf der Grundlage der aufgenommenen Spektren wird eine Energiekalibrierung des Systems durchgeführt. Dabei werden die bekannten Gammaenergien der Quellen mit den Kanalnummern verknüpft, die den Photopeaks entsprechen. Mittels linearer Regression wird eine Kalibriergerade bestimmt, die die Umrechnung von Kanalnummern in Energiewerte ermöglicht.\\

Nach erfolgreichem Abschluss der Kalibrierung wird die Halbwertsbreite (FWHM - Full Width at Half Maximum) der dominanten Peaks bestimmt. Dies umfasst zwei Linien aus dem $^{60}Co$-Spektrum, die 661-keV-Linie von $^{137}Cs$ sowie bestimmte, gut aufgelöste Linien aus dem Eu-152-Spektrum. Aus der Halbwertsbreite und der Peakenergie wird die Energieauflösung des Detektors berechnet und die Energieabhängigkeit der Auflösung untersucht. Als Nächstes wird das Peak-zu-Gesamt-Verhältnis für die Spektren von $^{137}Cs$ und $^{60}Co$ ermittelt. Hierzu wird die Fläche des Photopeaks ins Verhältnis zur Gesamtzahl der registrierten Ereignisse gesetzt. Dieses Verhältnis gibt Aufschluss über die Reinheit des Spektrums sowie über den relativen Anteil des Compton-Untergrunds und der Streuprozesse innerhalb des Detektors. Schließlich wird die Effizienz des NaI-Detektors mithilfe einer $^{137}Cs$-Quelle bestimmt. Die absolute Nachweiseffizienz wird unter Berücksichtigung der bekannten Aktivität, der Emissionswahrscheinlichkeit der 661-keV-Linie, der Messdauer sowie der Geometrie berechnet. Dies ermöglicht eine Abschätzung der Anzahl der emittierten Photonen, die tatsächlich vom Detektor registriert werden.

\subsection{Justage des HPGe-Halbleiterdetektor}

Zunächst wird das Signal des HPGe-Detektors direkt an der Vorverstärkerstufe mithilfe eines Oszilloskops analysiert. Dabei werden Signalform, Impulsdauer und Amplitude erfasst. Die Signalform wird durch die Bildung und Ansammlung von Elektron-Loch-Paaren im verarmten Germaniumkristall sowie durch die Ladungsansammlung im Vorverstärker bestimmt. Die Ausgangsspannung wird anschließend vom Hauptverstärker verarbeitet und in den Vielkanalanalysator eingespeist, um das Spektrum zu erhalten.\\

Im nächsten Schritt wurden Gammaspektren für die radioaktiven Elemente $^{60}Co$, $^{137}Cs$ und $^{152}Eu$ gemessen. Zusätzlich wurde unter Verwendung derselben Geometrie das Untergrundspektrum aufgezeichnet, um den Beitrag der Umgebungsstrahlung zu minimieren. Die Messdauer wurde so gewählt, dass scharfe, zuverlässige Photopeaks gewährleistet waren.\\

Auf der Grundlage dieser Spektren wird eine Intensitätskalibrierung des HPGe-Detektors durchgeführt. Dieser Prozess umfasst die Zuordnung der bekannten Gammaintensität zur Kanalnummer des Photopeaks und die Definition einer Umrechnungsfunktion mittels linearer Regression, wodurch die Umwandlung von Kanalnummern in Intensitätswerte ermöglicht wird.\\

Anschließend wird die Halbwertsbreite (FWHM - Full Width at Half Maximum) der Hauptkomponenten berechnet. Zu diesen Komponenten gehören zwei Peaks aus dem $^{60}Co$-Spektrum, der 661-keV-Peak von $^{137}Cs$ sowie mehrere intensive, gut definierte Peaks aus dem $^{152}Eu$-Spektrum. Unter Verwendung der FWHM und der Peakintensität werden Intensitätswerte berechnet, was die Bestimmung der optimalen Empfindlichkeit des HPGe-Detektors ermöglicht.\\

Das Peak-zu-Gesamt-Verhältnis für die $^{137}Cs$- und $^{60}Co$-Spektren wurde grafisch dargestellt. Hierfür werden die Positionen der Photopeaks relativ zur Anzahl der registrierten Ereignisse (Counts) normiert. Diese Auswertung erlaubt eine Beurteilung der spektralen Reinheit sowie der Größe des Compton-Untergrunds bei einem Halbleiterdetektor.\\

Der nächste Schritt umfasst die Bestimmung der Betriebsleistung des HPGe-Detektors. Hierfür wurden Spektren von $^{137}Cs$ und $^{152}Eu$ herangezogen. Unter Berücksichtigung der bekannten Aktivität, des Temperaturbereichs der jeweiligen Phase, der Messdauer und der Geometrie wurde das Effizienzverhältnis für verschiedene Energien berechnet und als Funktion der Gammaintensität ausgedrückt.\\

Schließlich wird der HPGe-Detektor zur Analyse einer unbekannten Probe, beispielsweise einer Bodenprobe, eingesetzt. Zu diesem Zweck wurde ein kontinuierlicher 12-stündiger Test unter Verwendung einer Bleiabschirmung durchgeführt, um messbare Ergebnisse zu erzielen. Die Untergrundkorrektur erfolgte auf der Basis einer einzelnen Nullmessung, die ohne Signalaufnahme durchgeführt wurde. Nach Abzug des Untergrundspektrums wurden die im Spektrum sichtbaren Gammalinien identifiziert und spezifischen Nukliden sowie deren zugehörigen Zerfallsreihen zugeordnet. Dies ermöglicht die Bestimmung, ob natürliche oder künstliche Radioaktivität in der Probe vorliegt.\\